\hypertarget{cryptographic-services}{%
\chapter{10. Cryptographic services}\label{cryptographic-services}}

The local-services chapter brokered the agent sockets a workstation
runs, and set one of them aside. An ssh-agent could not be handed
through the same facade as a keyring or a display socket, and the reason
it could not is the subject here. Signing is the capability where the
broker pattern meets its edge: for one kind of signing the edge is a
construction problem the framework solves, and for another it is a line
the framework will not cross. The difference between the two is the
whole of the chapter, and it is the difference between proving who you
are and vouching that something is so.

\hypertarget{the-signing-oracle}{%
\section{10.1 The signing oracle}\label{the-signing-oracle}}

The obvious way to let a kennel use the operator's SSH key is to expose
the agent socket into it, and it is wrong in a way that is worth being
precise about. The ssh-agent protocol carries an opaque blob to be
signed, not a destination. An agent, or a filtering broker placed in
front of one, can bound which key signs but never which host the
signature authenticates to, because the host is nowhere in what the
agent sees. A workload that holds the socket ignores the curated client
config it was given, opens the socket itself, and asks for a signature
over a challenge it built for a host of its own choosing. The key that
was meant for one server now authenticates the workload at every server
that accepts it. The agent is a destination-blind signing oracle, and
filtering the client the workload is free to bypass constrains nothing.
The conclusion the design draws is flat: the workload holds no key and
no agent socket, and it does not choose its own destination.

\hypertarget{re-origination}{%
\section{10.2 Re-origination}\label{re-origination}}

What replaces the agent is a bastion that re-originates the connection,
and the mechanism turns on a synthetic key that is itself the
destination. For each destination a kennel is granted, the toolchain
mints a throwaway ed25519 keypair at compile time; the private half
lands in the kennel's constructed \texttt{\textasciitilde{}/.ssh}, and
the public half is recorded in the signed policy and bound to a forced
command that already names the destination and the operator's host-side
options for reaching it. The binding is one \texttt{authorized\_keys}
line:

\begin{kennelcode}
restrict,pty,command="ssh '-i' '~/.ssh/id_gh' -- 'git@github.com' \"$SSH_ORIGINAL_COMMAND\"" ssh-ed25519 AAAA...
\end{kennelcode}

The destination is the literal \texttt{git@github.com} baked into the
command, fixed by which synthetic key authenticated and never parsed
from anything the workload sends. The workload controls only
\texttt{\$SSH\_ORIGINAL\_COMMAND}, the remote command, forwarded as a
single argument; it cannot rewrite the line, reach any other host, or
authenticate at all with a key the policy did not mint. There is no
oracle, because each key is a capability for exactly one destination.

That line lives in no file the bastion could read or a workload could
rewrite. On every authentication the bastion's \texttt{sshd} hands the
offered key to a root-owned helper and takes the answer as the key's
authorisation:

\begin{kennelcode}
AuthorizedKeysFile none
AuthorizedKeysCommand /opt/kennel/bin/kennel-akc %t %k
AuthorizedKeysCommandUser root
ExposeAuthInfo yes
\end{kennelcode}

The helper, \texttt{kennel-akc}, forwards the offered key to the running
\texttt{kenneld} over the control socket and prints back the
forced-command line bound to it, or nothing, in which case the login is
refused. The bindings exist only in the daemon's memory, sourced from
signed policy, so the answer is the daemon's live verified edge state
rather than a file the unprivileged bastion user could edit, and any
failure at all, no daemon or a garbled key or an unexpected reply,
prints nothing and fails the login closed. The helper is installed
root-owned because OpenSSH refuses an \texttt{AuthorizedKeysCommand}
that is not. The helper read the offered key from the \texttt{\%t\ \%k}
arguments; \texttt{ExposeAuthInfo\ yes} separately hands the forced
command the identity of the key that authenticated, as
\texttt{\$SSH\_USER\_AUTH}, so the lookup and the running session learn
the key by different routes, argv for the one and the environment for
the other.

The forced command is the only thing a session may do, and the rest of
the bastion's configuration is hardened to keep it that way:

\begin{kennelcode}
PubkeyAuthentication yes
PasswordAuthentication no
PermitRootLogin no
UsePAM no
AllowTcpForwarding no
X11Forwarding no
AllowAgentForwarding no
PermitTunnel no
PermitOpen none
AllowStreamLocalForwarding no
Subsystem sftp /bin/false
\end{kennelcode}

Publickey login only, no agent to forward and so none to re-export as an
oracle, every forwarding and tunnelling avenue shut, and SFTP wired to a
binary that does nothing. A session that authenticates runs its one
forced command and has no other move. The file is the surfaced template,
admin-tunable through a root-owned cascade but never weakenable by the
workload, which cannot write the paths it resolves from.

The kennel reaches the bastion over the same egress proxy chapter 8
builds, with no SSH-specific transport of its own: the kennel's network
policy allows only the bastion's loopback endpoint as a host service,
and \texttt{host-netproxy} dials it as it dials any other egress target.
On a connection the bastion's forced command runs, as the operator and
with no added privilege, an \texttt{ssh} to the bound destination,
dialed host-side and so outside the kennel's egress entirely, and that
outbound \texttt{ssh} signs with the operator's real key from the
operator's own host-side store, against the destination the bastion
fixed. The real key never enters the kennel. The kennel's
\texttt{\textasciitilde{}/.ssh} is a tmpfs subtree holding only the
disposable synthetic keys and a generated read-only config that points
every host at the bastion; the operator's real keys, real config, and
real \texttt{known\_hosts} are absent by construction rather than by a
deny rule, never built into the view (\texttt{construction-by-absence}).
Because the real signature is bound host-side to a destination the
workload could not choose, reusing one host's key against another is not
a rule that can be broken but a path that does not exist (T1.6).

\hypertarget{where-re-origination-stops}{%
\section{10.3 Where re-origination
stops}\label{where-re-origination-stops}}

Re-origination works for SSH because a signature there authenticates to
a host, and a host is a property the bastion can verify and bind a key
to. Commit signing has no such property, and the same fix does not reach
it. A signature on a commit is not aimed at a destination that could be
fixed; its worth is that the operator's identity vouches for the bytes,
and it incorporates itself into the hash it covers, so there is nothing
host-side a facade could check and bind its decision to. Behind the
technical mismatch is the axiom that settles the matter. A kennel exists
to contain code the operator does not trust. A signature made as the
operator, by code running inside that boundary, stamps the operator's
verified identity onto whatever the workload presents, a release, a
forged commit, malware, and it does so permanently, not for one live
session. That is a trust root placed inside the very boundary the
project exists to confine, which is incoherent: a claim with nothing
behind it.

So the framework does not build a signing facade, and it refuses to
pretend one is safe. An exposed gpg-agent is the same destination-blind
oracle as the ssh-agent, worse for stamping a durable artefact rather
than authenticating a transient session. A grant of it compiles and
runs, unlike the bus services chapter 9 turns away at compile, and the
split is the same rule §9.4 drew: the raw agent socket is a dumb connect
the operator owns, not a facade the framework builds and vouches for, so
it warns where a mediated signing service would be refused. The compiler
derives a risk from it that \texttt{kennel\ policy\ risks} surfaces
against the operator's reason: a loud footgun the operator may still
pull, not a sanctioned path (\texttt{footgun-warn-dont-forbid}). What
the axiom forbids is attesting as the operator, not signing as such, and
an operator who does want a kennel to sign has a clean way to get it.
The answer is a separate identity: an ordinary signing keypair the
operator generates for the workload, distinct from their own, its
private half granted through the ordinary read allowlist as a readable
file in the view. No new machinery and no agent, so no oracle to abuse;
the workload reads the key and signs in-process as that identity, and
the signature is worth exactly the trust of an identity the operator
stood up on purpose, read by a verifier as the workload's and never as
the operator's. The line the axiom draws runs between those two
identities, not against signing itself, which is the design volume's
rule (§8.2) in mechanism: refuse the standing to vouch as the operator,
grant without ceremony the standing to vouch as an identity the operator
raised for the purpose. Where no separate identity is wanted, the plain
path remains: the workload commits unsigned, and the human signs on
review, before the push, from outside the kennel.

\hypertarget{the-line-generalised}{%
\section{10.4 The line, generalised}\label{the-line-generalised}}

The split under SSH and signing is the general shape of every capability
the monitor brokers, and naming it keeps future facades honest. A
capability is authentication-shaped when it is a constrained,
host-verifiable act the monitor can bind to a checkable property: may I
reach this host, may I open this connection, may I render this frame.
Those are what interposition exists to carry, and the bastion is the
worked proof that a dangerous one can be made safe by binding it to the
property that makes it checkable. A capability is attestation-shaped
when its worth derives entirely from the trust of its origin: vouching,
signing, issuing a secret, trust that this is so. A kennel's origin is
confined and untrusted by construction, so a kennel that attests to
others is a trust root in the wrong place, and the category is refused
rather than merely unbuilt (\texttt{authentication-never-attestation}).
No broker the monitor grows, present or future, may be a secrets broker
or a signing service, and routing the act through a keyring or a TPM or
a vault does not rescue it, because it only moves the attestation to a
claim of being authorised to retrieve on the workload's behalf, which is
authority the broker still cannot vouch for.

The positive form of the same line is where trust material legitimately
enters a kennel. What a workload needs to hold, it holds because the
operator placed it there: a secret a kennel needs arrives as a signed
construction parameter, present in the view because construction put it
there, never handed to the kennel by a peer at runtime. A hardware token
follows the same rule from the other side. Its touch is the operator's
presence, so it is used host-side by the bastion's outbound
\texttt{ssh}, where the touch happens on the operator's own hardware,
and the device is never bound into the kennel, because a workload that
could drive the authenticator directly would have turned the presence
check into one more thing it signs.
